% =============================================================================
% Artigo Científico — Trabalho 02 FAED/PPGIA/PUC-PR
% Formato: IEEE duas colunas, máximo 6 páginas
% Template: IEEEtran (versão 1.8b)
% =============================================================================
\documentclass[conference]{IEEEtran}

\usepackage[utf8]{inputenc}
\usepackage[T1]{fontenc}
\usepackage[brazil]{babel}
\usepackage{cite}
\usepackage{amsmath,amssymb}
\usepackage{graphicx}
\usepackage{booktabs}
\usepackage{hyperref}
\usepackage{xcolor}
\usepackage{algorithm}
\usepackage{algorithmic}

\graphicspath{{figures/}}

% =============================================================================
\begin{document}

\title{Análise Comparativa de Algoritmos de Busca em Grafos Aplicados a
       Redes de Backhaul Sem Fio para Smart Grid}

\author{
  \IEEEauthorblockN{[Nome dos Autores]}
  \IEEEauthorblockA{
    Programa de Pós-Graduação em Informática Aplicada (PPGIA)\\
    Pontifícia Universidade Católica do Paraná (PUC-PR)\\
    Curitiba, Brasil\\
    \{email\}@ppgia.pucpr.br
  }
}

\maketitle

% =============================================================================
\begin{abstract}
Este artigo apresenta uma análise comparativa de cinco algoritmos de busca
em grafos — Dijkstra, Busca Gananciosa, A*, BFS e DFS — aplicados a uma rede
de backhaul sem fio para smart grid. A topologia experimental consiste em
25 nós distribuídos em três camadas (Access Points, Store-and-Forward e
Remotes), com pesos de aresta derivados de um custo composto baseado em
latência, throughput e taxa de perda de pacotes, parametrizado segundo
as especificações públicas do rádio GE MDS Orbit MCR 900~MHz e o modelo
de propagação Friis com sombreamento log-normal. Cinco pares
origem-destino foram avaliados em $\geq 5$ execuções independentes cada,
mensurando tempo de execução, nós expandidos e memória de pico. Os
resultados demonstram que A* reduz a expansão de nós em até [X]\% em
relação ao Dijkstra, mantendo otimalidade de solução, enquanto BFS e DFS
produzem caminhos sistematicamente subótimos em grafos ponderados.
\end{abstract}

\begin{IEEEkeywords}
algoritmos de busca, grafos ponderados, backhaul sem fio, smart grid,
Dijkstra, A*, heurística admissível
\end{IEEEkeywords}

% =============================================================================
\section{Introdução}

Redes de backhaul sem fio para smart grid enfrentam requisitos conflitantes
de latência, confiabilidade e custo operacional. A seleção de rotas nessas
redes pode ser modelada como um problema de busca em grafo ponderado, onde
os pesos dos enlaces refletem métricas físicas da camada de rádio
frequência (RF): RSSI, SNR, latência e taxa de perda de pacotes
\cite{rappaport2002}.

Algoritmos de busca em grafos diferem fundamentalmente quanto à utilização
de informação de custo (estratégias cegas vs. informadas) e à garantia de
otimalidade da solução \cite{russell2020}. A comparação empírica dessas
diferenças em um domínio de aplicação real — em vez de grafos sintéticos
genéricos — permite uma análise crítica mais fundamentada dos trade-offs
entre eficiência computacional e qualidade de solução.

Este trabalho implementa e compara Dijkstra \cite{dijkstra1959}, A*
\cite{hart1968}, Busca Gananciosa, BFS e DFS em Python puro (sem uso de
bibliotecas que forneçam a lógica dos algoritmos), sobre uma topologia de
25 nós representando uma rede de backhaul de smart grid com dados sintéticos
compatíveis com LGPD. As contribuições são:

\begin{itemize}
  \item Função de custo composto parametrizada para enlaces RF de 900~MHz;
  \item Heurística euclidiana admissível e consistente derivada dos dados;
  \item Avaliação estatística com $\geq 5$ rodadas por par de teste.
\end{itemize}

% =============================================================================
\section{Trabalhos Relacionados}

% TODO: referenciar artigos sobre roteamento em redes de smart grid e
% comparações de algoritmos de busca em grafos de comunicação.

% =============================================================================
\section{Modelagem da Rede e Função de Custo}

\subsection{Topologia}

A rede é modelada como um grafo não-direcionado $G = (V, E)$ com $|V| = 25$
nós e $|E| = 32$ arestas. Os nós pertencem a três categorias:

\begin{itemize}
  \item \textbf{AP} (3 nós): Access Points com conexão ao backbone cabeado;
  \item \textbf{SAF} (7 nós): repetidores Store-and-Forward;
  \item \textbf{Remote} (15 nós): concentradores de medidores inteligentes.
\end{itemize}

\subsection{Modelo de Propagação}

Os parâmetros de enlace são gerados pelo modelo de Friis para perda no
espaço livre (ITU-R P.525) com sombreamento log-normal:

\begin{equation}
  PL(d) = 20\log_{10}(d) + 20\log_{10}(f) + 20\log_{10}\!\left(\frac{4\pi}{c}\right) + X_\sigma
\end{equation}

onde $d$ é a distância em metros, $f = 928\,\text{MHz}$,
$X_\sigma \sim \mathcal{N}(0, 7^2)$~dB.

O throughput efetivo é estimado pela capacidade de Shannon com eficiência
prática de 55\%:

\begin{equation}
  C = 0{,}55 \cdot B \cdot \log_2(1 + \text{SNR}_\text{lin})
\end{equation}

com $B = 200\,\text{kHz}$ (modo narrow-band SCADA do GE MDS MCR
\cite{ge_mds_mcr}).

\subsection{Função de Custo Composto}

O peso de cada aresta combina três métricas de qualidade do enlace:

\begin{equation}
  w(u,v) = \alpha \cdot \ell_{ms} + \beta \cdot (1 - \tau_{norm}) \cdot 100 + \gamma \cdot \rho
  \label{eq:peso}
\end{equation}

onde $\ell_{ms}$ é a latência em ms, $\tau_{norm} = C / C_{max} \in [0,1]$
é o throughput normalizado, $\rho$ é a taxa de perda de pacotes (\%), e
$\alpha = 0{,}4$, $\beta = 0{,}4$, $\gamma = 0{,}2$.

% =============================================================================
\section{Algoritmos de Busca}

\subsection{Dijkstra}

Expande o nó de menor custo acumulado $g(n)$.
Complexidade: $O((V+E)\log V)$.
Garante solução ótima em grafos com pesos não-negativos \cite{cormen2022}.

\subsection{A*}

Expande o nó de menor $f(n) = g(n) + h(n)$.
A heurística euclidiana:

\begin{equation}
  h(n) = d_E(n, t) \cdot \min_{(u,v) \in E} \frac{w(u,v)}{d_E(u,v)}
  \label{eq:heuristica}
\end{equation}

é admissível ($h(n) \leq h^*(n)$) e consistente, garantindo otimalidade
sem re-expansão de nós \cite{hart1968}.

\subsection{Busca Gananciosa}

Expande apenas por $h(n)$, ignorando $g(n)$.
Rápida, mas não garante otimalidade.

\subsection{BFS e DFS}

BFS (fila FIFO) garante mínimo de saltos em grafos não-ponderados.
DFS (pilha LIFO) não oferece garantia alguma de otimalidade.
Ambos têm complexidade $O(V+E)$.

% =============================================================================
\section{Experimentos e Resultados}

\subsection{Configuração Experimental}

Cinco pares origem-destino cobrem cenários distintos do backhaul:

\begin{table}[h]
\caption{Pares de teste}
\label{tab:pares}
\centering
\begin{tabular}{lll}
\toprule
\# & Origem & Destino \\
\midrule
1 & REM-01 & AP-01   \\
2 & REM-14 & AP-02   \\
3 & REM-01 & REM-15  \\
4 & REM-07 & REM-12  \\
5 & REM-02 & REM-09  \\
\bottomrule
\end{tabular}
\end{table}

Cada par foi executado $N = 5$ vezes independentes por algoritmo.
Métricas coletadas: tempo de execução (ms), nós expandidos, custo do caminho,
comprimento (saltos) e memória de pico (KB).

\subsection{Resultados}

% TODO: substituir pela tabela gerada a partir de results/metrics/sumario.csv

\begin{table}[h]
\caption{Resultados médios — Par 1 (REM-01 $\to$ AP-01)}
\label{tab:resultados}
\centering
\begin{tabular}{lrrrr}
\toprule
Algoritmo & Tempo (ms) & Nós & Custo & Saltos \\
\midrule
Dijkstra   & -- & -- & -- & -- \\
A*         & -- & -- & -- & -- \\
Gananciosa & -- & -- & -- & -- \\
BFS        & -- & -- & -- & -- \\
DFS        & -- & -- & -- & -- \\
\bottomrule
\end{tabular}
\end{table}

% TODO: inserir figura gerada pelo Pyvis / matplotlib
\begin{figure}[h]
  \centering
  \includegraphics[width=\columnwidth]{topologia_backhaul}
  \caption{Topologia da rede de backhaul (25 nós). AP: vermelho; SAF:
           laranja; Remote: verde. Espessura das arestas proporcional ao custo.}
  \label{fig:topo}
\end{figure}

% =============================================================================
\section{Análise Crítica}

\subsection{Qualidade da Solução}

Dijkstra e A* retornam, por construção, o caminho de custo mínimo segundo
a função \eqref{eq:peso}. BFS e DFS não incorporam pesos na decisão de
expansão, produzindo caminhos com custo sistematicamente superior ao ótimo
nesta topologia.

\subsection{Eficiência Computacional}

A heurística \eqref{eq:heuristica} direciona A* para o destino, reduzindo
o número de nós expandidos em relação ao Dijkstra. Em grafos hierárquicos
como o backhaul modelado, onde a topologia reflete proximidade geográfica,
espera-se uma redução significativa de nós expandidos.

\subsection{Complexidade Assintótica}

\begin{table}[h]
\caption{Complexidade assintótica}
\label{tab:bigO}
\centering
\begin{tabular}{lcc}
\toprule
Algoritmo   & Tempo           & Espaço  \\
\midrule
Dijkstra    & $O((V+E)\log V)$ & $O(V)$ \\
A*          & $O((V+E)\log V)$ & $O(V)$ \\
Gananciosa  & $O((V+E)\log V)$ & $O(V)$ \\
BFS         & $O(V+E)$        & $O(V)$  \\
DFS         & $O(V+E)$        & $O(V)$  \\
\bottomrule
\end{tabular}
\end{table}

\subsection{Trade-offs para Backhaul Smart Grid}

Em redes de smart grid, o custo de falha de entrega de um telegrama SCADA
(controle) supera o custo computacional de qualquer algoritmo.
Portanto, a prioridade é otimalidade e confiabilidade do caminho.
A* apresenta o melhor equilíbrio: garante otimalidade com custo
computacional reduzido pela heurística geográfica.

% =============================================================================
\section{Conclusão}

% TODO

% =============================================================================
\bibliographystyle{IEEEtran}
\bibliography{references}

\end{document}
